% !TeX root = ../main.tex

\chapter{项目概述}
\label{sec:project-overview}

\section{项目背景}

赛题聚焦大模型驱动的个性化资源生成与学习多智能体系统，面向高等教育中真实而持续的学习需求。随着专业课程资源不断丰富，学生能够接触到教材、课程视频、习题、病例材料和学术文献，却也需要自行完成资料筛选、学习顺序安排、理解检验和复习节奏管理。学习过程中产生的目标、疑问、作答和进度常常分散在不同平台，前后学习活动缺少连续衔接，学生很难及时获得与自身阶段相适应的下一步行动建议。

口腔医学的学习任务进一步放大了这一困难。课程体系横跨基础医学、临床诊疗、材料科学和预防保健等多个方向，多门核心学科之间存在紧密关联。一个疾病主题往往牵连组织解剖、病因机制、临床表现、检查依据、鉴别思路、治疗原则及预后维护等多层知识。学生在备考、专题复习和病例研读时，既要掌握清晰的概念边界，又要理解知识之间的关联。当学习内容只能以孤立的段落、题目和视频呈现时，知识体系容易停留在记忆层面，难以形成可迁移的专业理解。

近年来，生成式人工智能已进入医学与口腔医学教育的案例讨论、问题导向学习和苏格拉底式追问等场景，为知识解释、互动讨论和学习资源组织带来新的空间\cite{wei2025aipbl,liu2026socratic}。面向口腔健康教育的大语言模型研究也指出，这类工具能够支持信息获取、个性化解释与学习互动，同时需要认真对待内容准确性、教学适配性和使用边界\cite{sharmin2025oralhealthllm}。牙科教育题目表现的元分析进一步表明，模型表现会随着模型版本、题目类型和题源变化而波动\cite{doubleday2026llmaccuracy}。这些研究提示，专业学习系统需要将生成能力放入明确的课程范围和专业资料环境中，让学习者获得可跟随的学习路径与持续反馈。

传统学习资源通常按照资料类型组织，教材提供系统知识，题库用于检验，课程视频承担讲解，文献服务专题查阅。学生需要在多个资源入口之间反复切换，自行判断当前应该学习什么、优先复习什么，以及一道错题反映了哪些知识缺口。通用问答工具能够回应即时问题，却难以持续理解学习者正在学习的课程位置、已掌握内容和近期安排。对于节奏差异显著的专业学习者，真正需要的是能够将知识内容、学习行为和后续任务自然串联起来的学习环境。

OralSage 以口腔医学为垂直场景，围绕“确定目标、进入课程、理解知识、检验掌握、回顾薄弱内容、安排后续学习”构建连续服务。项目把权威资料、结构化知识、课程资源和学习记录汇聚到统一工作台，让学习者能够从一个具体主题出发，在同一学习过程中完成阅读、听讲、提问、练习和复盘。这样的场景选择既回应了赛题对个性化资源生成与多智能体协同的要求，也面向口腔医学学习中知识密集、路径差异明显和复习任务持续变化的实际需求。

\section{项目实现}

OralSage 构建了以学习者为中心的口腔医学智能学习平台。学习者可以从专题输入、课程章节或日历任务进入系统，系统结合当前目标、可用时间、学习偏好、已学进度和掌握情况组织学习内容。一个主题进入课程后会形成清晰的学习单元和学习顺序，首个单元可以直接开始学习。讲义、课件、讲稿、音频、思维导图和练习围绕同一主题展开，学习者能够根据当下的时间与习惯选择合适的学习方式，并始终保留课程主线。

平台采用多智能体协作组织连续的学习活动。课程规划、资料组织、分节讲解、课程内答疑和理解检验由不同学习型 Agent 分别承担，并在同一学习会话中相互接续。学习者提出问题时，系统围绕当前章节补充解释，完成答题后可以看到正误反馈、错因分析和复习建议，完成一个单元后能够顺畅进入下一阶段。多种服务共享课程位置和学习记录，使学习者无需反复描述主题与进度，能够把注意力放在知识理解和学习推进上。

个性化能力建立在持续更新的学习画像之上。初始测评帮助系统了解学习目标、每日可用时间、知识进入方式、讲解深度、练习倾向和计划支持需求。后续的课程学习、提问、作答、复习和计划执行逐渐沉淀为学习状态，反映学习者的掌握情况、易错内容、学习偏好和执行节奏。学习计划据此将继续课程、重点回顾、定向练习和阶段自测组织为明确任务，学习状态页面则通过进度、趋势、能力雷达、学习活动和时间线呈现个人学习轨迹，让抽象的学习变化变得可见、可跟随。

针对临床沟通与口头表达难以在普通学习平台中反复练习的问题，OralSage 增设独立的数字问诊页面。讯飞实时数字人依据脱敏病例扮演患者，患者 Agent 理解学生的语音或文字追问并按病例事实回答，识别到正确诊断后由判官 Agent 接续病理机制、诊断依据、鉴别诊断、治疗原则和风险随访考核。分项评分、遗漏要点和相关知识节点会进入学习画像，后续病例选择继续参考这些训练结果，使临床训练与个性化学习形成持续联系。

专业知识服务为上述学习过程提供稳定支撑。当前知识底座覆盖 14 门口腔医学学科，汇集 9,907 个知识概念、5,758 条知识关系、21,737 个专业知识片段和 3,161 道关联题目。系统通过知识图谱组织课程主题与概念联系，通过多路知识检索向课程规划、学习资料、答疑和练习反馈提供相关内容，并保留资料来源信息，方便学习者在学习过程中进一步查阅。自研 ClawBot 将网页工作台延展至微信学习场景，支持账户绑定和学习计划推送，帮助学习者在日常沟通环境中接收任务提醒并保持学习节奏。

在专业视觉学习方面，团队围绕口腔全景片等牙科影像的宽幅构图、细节密集和目标细微等特点，构建了 OralPath-7B 口腔影像视觉语言模型。模型以 Qwen2.5-VL-7B-Instruct 为基础完成动态分辨率 LoRA 微调，训练资料来自两万余张高质量标注牙科影像，覆盖 31 类牙科结构与常见病变目标，并重构为五万余条图像、任务指令和标准回答相结合的图文样本。模型能够将牙体、牙周组织、根尖区、牙槽骨、阻生牙、修复体和常见病变等视觉信息组织为目标定位与中文影像摘要，为影像学习页面提供由整体观察进入局部辨识和专业讨论的学习支持。该能力服务于口腔医学教学与资料分析，帮助学习者建立规范的影像阅读思路。

围绕专题学习和专业能力拓展，平台还提供影像学习辅助、病理问答、病例资料分析、文献阅读、知识浏览、统计查询和病历文书等工具 Agent。它们服务于资料查阅、专题讨论和专业表达训练，使学习者能够在课程学习之外继续探索相关知识。所有临床相关内容均定位于教学学习与资料参考，帮助学生建立专业思考框架。

\section{核心功能概览}

OralSage 将课程学习、学习管理、专业资料服务和跨端陪伴组织为统一产品体验。各项核心功能及其面向学习者的价值如表\ref{tab:core-function-overview}所示。

{\small
\renewcommand{\arraystretch}{1.18}
\begin{longtable}{p{0.23\linewidth}p{0.30\linewidth}p{0.37\linewidth}}
  \caption{OralSage 核心功能概览}\label{tab:core-function-overview}\\
  \toprule
  功能类别 & 核心能力 & 学习者可获得的体验与价值 \\
  \midrule
  \endfirsthead
  \multicolumn{3}{c}{续表\ \thetable}\\
  \toprule
  功能类别 & 核心能力 & 学习者可获得的体验与价值 \\
  \midrule
  \endhead
  \bottomrule
  \endfoot
  学习型 Agent 群 & Planner Agent 组织课程主题、学习单元与学习顺序 & 复杂专题被拆解为具有递进关系的课程路径，学习者能够明确当前学习重点和后续方向。 \\
  学习型 Agent 群 & Resource Factory Agent 准备讲义、课件、讲稿、音频、思维导图和练习材料 & 多种学习资料围绕同一课程单元协同呈现，支持精读、浏览、听讲、回顾与主动检验等不同学习方式。 \\
  学习型 Agent 群 & Tutor Agent 提供分节讲解、重点提示和随堂互动 & 学习者能够沿课程节奏理解概念及其前后联系，在连续讲解中建立稳定的知识框架。 \\
  学习型 Agent 群 & QA Agent 处理当前课程中的即时追问 & 提问获得与正在学习的章节相衔接的补充说明，答疑结束后可以自然回到原有学习进度。 \\
  学习型 Agent 群 & Assessment Agent 完成练习反馈、错因分析和补救建议 & 一次作答能够转化为可理解的掌握反馈，帮助学习者定位需要回顾的知识内容。 \\
  数字问诊与临床训练 & 数字人患者、患者 Agent 和判官 Agent 组织语音或文字问诊及五域表达考核 & 学习者可以练习病史采集、诊断判断与完整口头表达，分项结果继续影响学习画像和后续病例选择。 \\
  学习画像 & 初始测评与动态学习记录共同描述目标、时间、偏好、掌握情况和薄弱内容 & 系统能够在后续课程、讲解、练习和任务安排中参考个人状态，形成更贴合学习节奏的服务。 \\
  学习计划与状态跟进 & 日历任务、课程进度、复习安排、执行记录和可视化学习状态 & 学习者可以查看本日任务、近期进展与待回顾内容，将阶段反馈转化为下一次可执行的学习行动。 \\
  自研 ClawBot & 网页工作台与微信学习场景的账户绑定和计划推送 & 学习计划能够进入日常沟通环境，帮助学习者及时接收提醒并持续完成学习安排。 \\
  知识库 & 覆盖 14 门口腔医学学科的教材、规范、手册、知识概念和关联题目 & 课程学习、知识查阅与练习反馈拥有统一的专业内容基础，学习者能够从主题资料进一步追溯相关知识。 \\
  知识图谱与知识检索 & 知识概念、关系、资料片段和题目之间的结构化关联及多路检索服务 & 学习者可以沿概念联系理解专题，也能快速定位与当前问题相关的教材资料和练习内容。 \\
  专业工具 Agent & 影像学习辅助、病理问答、病例资料分析与文献阅读 & 面向专题讨论、病例研读和前沿资料查阅提供专业学习支持，帮助学生拓展课程知识的应用视野。 \\
  专业工具 Agent & 知识浏览、统计查询与病历文书工具 & 面向知识探索、数据理解和专业表达训练提供补充服务，丰富口腔医学学习的实践化场景。 \\
  可视化交互 & 课程会话、知识图谱、学习画像、学习计划和专业工具工作台 & 抽象的课程结构、知识联系和个人进展以清晰页面呈现，学习者能够在不同学习任务之间保持顺畅切换。 \\
\end{longtable}
}

通过上述功能，OralSage 将一次主题学习延展为可持续推进的学习过程。学习者既可以按课程路径系统学习，也可以围绕眼前问题获得专业资料支持，并在学习反馈和计划提醒的陪伴下不断回到下一步行动。
% !TeX root = ../main.tex
